\subsection{NPU 推理与 RKNN 运行时}\label{sec:npu}

经过缩放与色彩转换的那一帧最终要送进检测网络，承担这一步的是 RK3588 片上集成的神经网络
加速器。它由三个独立的计算核组成，每个核运行在 1.0\,GHz，专门以低精度整数算术高吞吐地执行
卷积、矩阵乘与激活这类张量算子。直接在 CPU 上跑同一张网络会让感知与控制争抢通用核，一帧的
代价高到流水线无法承受，把这部分计算交给专用加速器，CPU 才能腾出来去做决策与运动控制。

加速器自身只认一种以低精度整数权重编排好的指令流，普通深度学习框架训练出的模型并不能直接喂
给它。瑞芯微为此提供了一套运行时与模型格式，运行时以用户态共享库 \texttt{librknnrt.so} 的形式
存在，模型则由 rknn-toolkit2 这一离线工具把常见框架的网络转换为 RKNN 格式，转换过程顺带把
浮点权重量化为 INT8 并记下每个张量的量化标度。应用通过这套运行时的接口加载模型、提交一帧、
取回结果，运行时再把这些请求翻译成对加速器的命令提交与内存操作，经设备节点下达到内核。

这条路径在本系统里落在 \texttt{/dev/card1} 这个设备节点上，它在内核中表现为一个 DRM（图形显示
子系统）字符设备。运行时与内核之间的交互全部以 ioctl 命令完成，其命令号按 DRM 的约定编码，
低于 \texttt{0x40} 的是 DRM 核心命令如查询驱动版本，落在 \texttt{0x40} 到 \texttt{0xA0} 之间的
才是 RKNPU 这个驱动自己的命令，节点据此把版本查询与厂商私有的内存与提交命令分派到不同的处理
路径。RKNPU 的私有命令共六类，分别负责内存对象的创建、映射、销毁、缓存同步，以及任务的动作
查询与提交，对应 \texttt{rknpu\_driver\_ioctl} 里的 \texttt{MemCreate}、\texttt{MemMap}、
\texttt{MemDestroy}、\texttt{MemSync}、\texttt{Action} 与 \texttt{Submit}。

加速器看到的输入张量并不是一段裸内存，而是经由图形子系统的缓冲对象（GEM，graphics 子系统里
统一管理显存与设备内存的句柄）来描述。运行时先以一次内存创建命令换得一个句柄，再以这个句柄
向内核请求把对应的物理内存映射进自己的地址空间，此后对这块张量的读写就直接落在加速器要访问的
那段内存上。这层缓冲对象正是让上一级的图像加速器与本级共享同一段内存的关键，输入缓冲以句柄
在两者之间传递，加速器与图像准备阶段读写的是同一块物理页而非各自的副本，整条流水线因此免去了
一次拷贝。

\paragraph{让每次推理都能接续前一次}
这个节点与其驱动在本系统里本已存在，但要真正承载一帧接一帧的连续推理，缓冲对象在两处的行为
必须被修正，否则总是上一帧能跑，下一帧就失败。第一处在映射时。运行时按句柄请求映射，节点把
句柄换算回缓冲对象的物理地址并交给上层去建立页表，节点据缓存策略决定这段映射是可缓存还是非
缓存，写合并因为内核尚未提供对应的页表项模式而被退化为非缓存处理，以免误把它当作可缓存而读到
陈旧数据。映射建立时若不正确处理这段区域的缺页，第一次访问就会触发无法收场的异常，使紧随其后
的那次推理取不到输入。第二处在销毁时。每个缓冲对象都带引用计数，运行时取回结果后会销毁这一帧
的临时对象，若销毁时计数处理有误，被提前回收或重复释放的那段内存会在下一帧被再次创建时拿到错
位的地址。把映射时的缺页处理与销毁时的引用计数这两处都修正之后，连续推理才真正稳定下来。\textbf{这部分
对 RKNPU 的 DRM 接入与 GEM 缓冲对象的修复以合并请求 \href{https://github.com/rcore-os/tgoskits/pull/1351}{\#1351}
与 \href{https://github.com/rcore-os/tgoskits/pull/1364}{\#1364} 提交至上游，均已合入。}

\paragraph{在内核侧为各阶段计时}
为了弄清一帧的时间究竟花在哪个 ioctl 阶段，节点里加进了一个可选的计时特性 \texttt{profile-rknpu}。
它在提交、内存创建、内存同步等命令的处理两端各取一次 \texttt{monotonic\_time\_nanos} 单调时钟，
算出该阶段的微秒耗时，并用一组 \texttt{AtomicUsize} 计数器以 \texttt{Relaxed} 次序无锁地为每条
记录编号，超过给定条数后便不再打印。整个计时通路不持有任何锁，因此既不改变被测路径的时序，
也不会在多核并发提交时相互阻塞，关掉这个特性时这些代码完全不参与编译，正常运行不付出任何代价。

\paragraph{在整数域筛除背景}
最显著的一处优化落在取回结果这一步，做在应用一侧。加速器算完后吐出的是 INT8 张量，要还原成
有物理意义的浮点置信度，需要按各张量的量化标度做一次反量化。一帧的检测头在三个尺度上铺开约
四十万个网格值，其中绝大多数是背景，没有任何目标。若让运行时在取回时把这四十万个值统统反量化
为浮点，绝大部分算力都花在了注定被丢弃的背景上。本系统把取回时的浮点化关掉，令运行时直接交回
原始的 INT8 张量，再把置信度阈值反向量化到整数域，在整数比较中先把低于阈值的网格全部筛除，
只对幸存下来的那几十个网格做反量化。反量化是单调映射，在整数域比较置信度的大小与在浮点域比较
完全等价，因此输出与逐点反量化的结果逐位一致。这样一来取回结果一步从约 \texttt{12}\,ms 降到约
\texttt{1}\,ms，而在全部 285 张测试图上，整数域路径得到的检测框与浮点路径完全相同。

这张网络以卷积类算子为主，约七成的加速器时间都花在卷积上。它的规模不大，单帧的负载基本压在
一个计算核上，再开第二、第三个核并不会带来收益，那两个核多数时候是空闲的，因此实际推理始终
只用到一个核，加速器全程保持在 1.0\,GHz。也正因为时钟固定在最高频不做动态降频，推理这一步
（\texttt{run}）在本系统上与 Linux 在同一时钟下处于同一水平，约 \texttt{4.7}\,ms，并不更快，
真正被压下来的是它两侧的取回与内存操作，而非加速器的计算本身。

\repfig{npu-rknn.png}{运行时经 \texttt{/dev/card1} 提交一帧的路径，输入张量以 GEM 缓冲对象在图像准备阶段与加速器之间零拷贝共享，取回结果在整数域筛除背景后只对幸存网格反量化}{fig:npu}
